\minitoc 

% ========================================================================================================
% Numérisation et Représentation 2D

\section{Numérisation et Représentation de l'Image 2D}

\subsection{Échantillonnage, Quantification et Modèle matriciel}

L'acquisition d'une image continue $i(x,y)$ passe par deux étapes de discrétisation :
\begin{enumerate}
    \item \textbf{Échantillonnage spatial} : mesure selon un pas régulier $(d_x, d_y)$ respectant la condition de Shannon : $i(m,n) = i(m d_x, n d_y)$.
    \item \textbf{Quantification d'amplitude} : projection de la valeur d'intensité continue sur un ensemble discret de $L$ niveaux de gris. Généralement, $L = 2^k$ (image codée sur $k$ bits).
\end{enumerate}

\begin{definition}[Image matricielle et Pixel]
    Une image numérique monochrome est une matrice $I(l,c)$ de $M$ lignes et $N$ colonnes :
    $$ l \in \llbracket 0, M-1 \rrbracket \quad (\text{axe vertical } y), \qquad c \in \llbracket 0, N-1 \rrbracket \quad (\text{axe horizontal } x) $$
    La cellule élémentaire $(p, I(p))$ avec $p = (l,c)$ est appelée \textbf{pixel}.
\end{definition}

\begin{proposition}[Poids mémoire brut]
    La taille brute $B$ en bits nécessaire pour stocker sans compression une image $M \times N$ sur $k$ bits est :
    \begin{equation}
        B = M \times N \times k \quad \text{bits}
    \end{equation}
    Pour une image carrée $N \times N$, $B = N^2 k$ bits.
\end{proposition}

\subsection{Typologie d'images}

\begin{itemize}
    \item \textbf{Image binaire} : $I(l,c) \in \{0, 1\}$ ($1$ bit par pixel).
    \item \textbf{Image à niveaux de gris (8 bits)} : $I(l,c) \in \llbracket 0, 255 \rrbracket$.
    \item \textbf{Image couleur (RGB - Synthèse additive)} : chaque pixel est un vecteur tridimensionnel :
        $$ I(l,c) = \begin{bmatrix} R(l,c) \\ G(l,c) \\ B(l,c) \end{bmatrix} \in \llbracket 0, 255 \rrbracket^3 $$
    \item \textbf{Luminance et Valeur moyenne} :
        \begin{align}
            \text{Luminance } L(l,c) &= 0.299 \, R(l,c) + 0.587 \, G(l,c) + 0.114 \, B(l,c) \\
            \text{Valeur } V(l,c) &= \frac{R(l,c) + G(l,c) + B(l,c)}{3}
        \end{align}
    \item \textbf{Fausses couleurs} : application d'une table de correspondance (LUT - \emph{Look-Up Table}) sur une image scalaire pour accentuer des contrastes à l'œil humain.
\end{itemize}

% ========================================================================================================
% Traitements ponctuels (Point-à-point)

\section{Traitements Ponctuels (Point-à-point)}

Un opérateur point-à-point transforme l'intensité d'un pixel indépendamment de ses voisins : $s = T(r)$, où $r$ est le niveau de gris d'entrée et $s$ le niveau de gris de sortie.

\begin{proposition}[Transformations usuelles]
    \begin{itemize}
        \item \textbf{Luminosité} :
            \begin{equation}
                s = \operatorname{tronc}(r + g), \quad g \in \Z
            \end{equation}
        \item \textbf{Contraste linéaire} :
            \begin{equation}
                s = \operatorname{tronc}(a [r - 127] + 127), \quad a > 1 \text{ (augmente)}, \quad 0 \leqslant a < 1 \text{ (diminue)}
            \end{equation}
        \item \textbf{Négatif} :
            \begin{equation}
                s = (L - 1) - r = 255 - r
            \end{equation}
        \item \textbf{Seuillage binaire} :
            \begin{equation}
                s = \begin{cases} 1 & \text{si } r \geqslant \text{seuil} \\ 0 & \text{si } r < \text{seuil} \end{cases}
            \end{equation}
        \item \textbf{Transformation logarithmique} :
            \begin{equation}
                s = c \log(1 + r) \quad (\text{dilatation des niveaux sombres, compression des niveaux clairs})
            \end{equation}
        \item \textbf{Loi de puissance (Correction Gamma)} :
            \begin{equation}
                s = c \, r^\gamma \quad (\gamma < 1 \text{ éclaircit l'image}, \quad \gamma > 1 \text{ l'assombrit})
            \end{equation}
    \end{itemize}
\end{proposition}

% ========================================================================================================
% Histogramme et Égalisation

\section{Histogramme et Égalisation statistique}

\subsection{Définitions statistiques}

\begin{definition}[Histogramme, ddp et Fonction de répartition]
    Pour une image de taille $M \times N$ comportant $L = 256$ niveaux de gris :
    \begin{itemize}
        \item \textbf{Histogramme} $h(g)$ : nombre de pixels d'intensité $g \in \llbracket 0, L-1 \rrbracket$.
        \item \textbf{Densité de probabilité (ddp / pdf)} :
            \begin{equation}
                p(g) = \frac{h(g)}{M \times N}, \quad \text{avec } \sum_{g=0}^{L-1} p(g) = 1
            \end{equation}
        \item \textbf{Fonction de répartition (FDR / CDF)} :
            \begin{equation}
                P(g) = \sum_{k=0}^g p(k) = \frac{1}{M \times N} \sum_{k=0}^g h(k)
            \end{equation}
            $P(g)$ est une fonction croissante telle que $P(0) = p(0)$ et $P(L-1) = 1$.
    \end{itemize}
\end{definition}

\subsection{Égalisation d'histogramme}

L'objectif est d'étaler au maximum la dynamique des niveaux de gris pour obtenir un histogramme plat (uniforme).

\begin{theorem}[Formule d'égalisation]
    L'égalisation d'histogramme utilise directement la fonction de répartition $P(g)$ comme LUT de transformation :
    \begin{equation}
        \boxed{J(l,c) = (L - 1) \cdot P_I[I(l,c)]}
    \end{equation}
    Dans le cas général où l'on souhaite recadrer l'intervalle d'entrée $[m_I, M_I]$ vers $[m_J, M_J]$ :
    \begin{equation}
        J(l,c) = (M_J - m_J) \frac{P_I[I(l,c)] - P_I(m_I)}{1 - P_I(m_I)} + m_J
    \end{equation}
\end{theorem}

% ========================================================================================================
% Filtrage spatial 2D

\section{Filtrage Spatial 2D}

\subsection{Convolution 2D et Traitement des bords}

Le filtrage spatial consiste à déplacer un masque (noyau) $w(s,t)$ de dimension $(2a+1) \times (2b+1)$ sur l'image :
\begin{equation}
    g(x,y) = \sum_{s=-a}^a \sum_{t=-b}^b w(s,t) f(x+s, y+t)
\end{equation}

\begin{prop}[Gestion des bordures de convolution]
    Lors du passage du masque sur les bords de l'image, les valeurs extérieures sont interpolées selon trois stratégies :
    \begin{itemize}
        \item \textbf{Zero-padding} : assignation de zéros en dehors de l'image (simple mais crée un artéfact de bordure sombre).
        \item \textbf{Extension périodique} : réplication de l'image par pavage toroïdal (en accord avec la TFD 2D).
        \item \textbf{Extension miroir (symétrie)} : réplication par réflexion par rapport au bord (atténue fortement les discontinuités de bordure).
    \end{itemize}
\end{prop}

\subsection{Filtres de lissage}

\begin{itemize}
    \item \textbf{Filtre moyenneur (linéaire)} : masque uniforme de coefficients normalisés $w = \frac{1}{(2K+1)^2} \mathbf{1}$. Pour une fenêtre $3 \times 3$ :
        $$ w = \frac{1}{9} \begin{bmatrix} 1 & 1 & 1 \\ 1 & 1 & 1 \\ 1 & 1 & 1 \end{bmatrix} $$
        Réduit le bruit blanc gaussien mais floute les contours.
    \item \textbf{Filtre médian (non-linéaire)} : remplace la valeur du pixel central par la médiane des valeurs triées du voisinage.
        \begin{equation}
            J(x,y) = \operatorname{médiane}\{ I(x+s, y+t), (s,t) \in \mathcal{V} \}
        \end{equation}
        Optimal pour éliminer le bruit impulsionnel dit « poivre et sel » tout en préservant la netteté des contours.
\end{itemize}

\subsection{Filtres de réhaussement de contours (Dérivées spatiales)}

Les contours correspondent à des variations rapides d'intensité, mises en évidence par dérivation spatiale.

\subsubsection{Opérateur Laplacien (Dérivée seconde)}

Le Laplacien continu $\nabla^2 f = \frac{\partial^2 f}{\partial x^2} + \frac{\partial^2 f}{\partial y^2}$ est discrétisé par :
\begin{equation}
    \nabla^2 f(x,y) = f(x+1, y) + f(x-1, y) + f(x, y+1) + f(x, y-1) - 4 f(x,y)
\end{equation}

\begin{proposition}[Masques du Laplacien et Réhaussement]
    \begin{itemize}
        \item Masque standard (4 voisins) et masque étendu (avec diagonales, 8 voisins) :
            $$
            \nabla^2_{4} = \begin{bmatrix} 0 & 1 & 0 \\ 1 & -4 & 1 \\ 0 & 1 & 0 \end{bmatrix}, \qquad 
            \nabla^2_{8} = \begin{bmatrix} 1 & 1 & 1 \\ 1 & -8 & 1 \\ 1 & 1 & 1 \end{bmatrix}
            $$
        \item \textbf{Image réhaussée} : pour restaurer la dynamique globale tout en accentuant les contours, on retranche le Laplacien :
            \begin{equation}
                g(x,y) = f(x,y) - \nabla^2 f(x,y)
            \end{equation}
            Ce qui donne en un masque unique condensé :
            $$
            w_{\text{réhauss}} = \begin{bmatrix} 0 & -1 & 0 \\ -1 & 5 & -1 \\ 0 & -1 & 0 \end{bmatrix}
            $$
        \item \textbf{Masquage flou (\emph{Unsharp Masking})} : $g(x,y) = A f(x,y) - \nabla^2 f(x,y)$ avec $A \geqslant 1$ pour amplifier le contraste.
    \end{itemize}
\end{proposition}

\subsubsection{Opérateur Gradient et Filtres de Sobel (Dérivée première)}

Le gradient est le vecteur $\nabla f = [\frac{\partial f}{\partial x}, \frac{\partial f}{\partial y}]^T = [G_x, G_y]^T$. Son amplitude est approchée en pratique par $\|\nabla f\| \approx |G_x| + |G_y|$.

\begin{definition}[Masques de Sobel]
    Le filtre de Sobel applique deux masques de convolution directionnels :
    \begin{equation}
        G_x = \begin{bmatrix} -1 & 0 & 1 \\ -2 & 0 & 2 \\ -1 & 0 & 1 \end{bmatrix}, \qquad
        G_y = \begin{bmatrix} -1 & -2 & -1 \\ 0 & 0 & 0 \\ 1 & 2 & 1 \end{bmatrix}
    \end{equation}
    $G_x$ détecte les contours verticaux et $G_y$ les contours horizontaux.
\end{definition}